
\section{Validation, Métrologie \& Traçabilité (Section 10)}

\subsection{Objectifs}
Fournir des garanties quantitatives sur la \textbf{validité} du pipeline (exactitude, fidélité, robustesse) et une \textbf{traçabilité complète} (provenance $\rightarrow$ décisions algorithmiques $\rightarrow$ exports), conforme aux pratiques labo/industrie (ICH Q2/USP/ISO).

\subsection{Plan d'expériences (DOE)}
\begin{itemize}[nosep]
\item \textbf{Répétabilité} (même jour, même instrument, même opérateur): $n\ge 10$ réplicats/échantillon.
\item \textbf{Reproductibilité intermédiaire} (jours différents/opérateurs différents).
\item \textbf{Inter-instruments} (sites/instruments différents, $\lambda_L$, objectifs, filtres).
\item \textbf{Stabilité} vs temps (contrôle dérive, recalibration périodique).
\end{itemize}

\subsection{Métriques de performance}
\begin{itemize}[nosep]
\item \textbf{Exactitude} (biais): moyenne $\hat{C}-C$ (Sec.~9) et test de Bland--Altman (LOA).
\item \textbf{Fidélité} (précision): écart-type intra-condition; \textbf{r\%} (\%RSD) cible $\le 5$--$10\%$ selon cas.
\item \textbf{Robustesse}: variation sous perturbations contrôlées (gain/offset, \emph{baseline} $\lambda$, SG-$M$, bruit).
\item \textbf{LOD/LOQ vérifiées}: expérience de dilutions en blanc (Sec.~9) avec taux de détection/quantification observés.
\item \textbf{Taux d'acceptation QA}: \% spectres \texttt{acceptés} (Sec.~3) et causes de rejet.
\item \textbf{Concordance identification}: macro-F1 (Sec.~8), ECE (calibration), TPR@FPR pour OOD.
\end{itemize}

\subsection{Gage R\&R (ANOVA)}
Pour une mesure $Y_{ijk}$ (pièce $i$, opérateur $j$, réplication $k$), ANOVA à effets aléatoires $\Rightarrow$ composantes de variance: \textit{repeatability} $\sigma^2_e$, \textit{reproducibility} opérateur $\sigma^2_o$, interaction pièce$\times$opérateur $\sigma^2_{po}$. Pourcentage de variation: $\mathrm{PV} = 100 \times \sigma_\text{gage}/\sigma_\text{total}$. Cible: PV $\le 10\%$ (bon), 10--30\% (marginal), $>30\%$ (à améliorer).

\subsection{Stabilité \& contrôle dérive}
\begin{itemize}[nosep]
\item \textbf{Shewhart/EWMA} sur $\mu_{520.7}$ (Si) et FWHM instrumentale: alarmes 1--2--3$\sigma$.
\item \textbf{Recalibration} (Sec.~4) si dérive $|\Delta \nu| > \tau_{\mathrm{drift}}$ persiste (ex. $>0.5$~\cm).
\item \textbf{Harmonisation} inter-instruments: correction de gain/offset et adaptation de domaine (Sec.~8).
\end{itemize}

\subsection{Robustesse (sensibilité aux hyperparamètres)}
\begin{itemize}[nosep]
\item Balayer $\lambda$ (AsLS), fenêtre SG, seuils QA, paramètres de fit Voigt; mesurer l'impact sur $A,\mu,\mathrm{FWHM}$ et sur les sorties Sec.~8/9.
\item \textbf{Critère} : variations $\Delta$ dans des \textit{bornes acceptables} (p.ex. $|\Delta \mu|<0.3$~\cm; $\Delta$F1$<2$~pts; $\Delta$RMSE$<5\%$).
\end{itemize}

\subsection{Traçabilité \& intégrité des données}
\begin{itemize}[nosep]
\item \textbf{Provenance} : \texttt{manifest.json} (hash SHA-256, versions, métas) et logs structurés (stade, config, seeds).
\item \textbf{Chaîne d'horodatage} : UTC, identifiants \texttt{run\_uuid}, \texttt{sample\_id}, \texttt{instrument\_id}.
\item \textbf{Reproductibilité} : \texttt{config.yaml} versionnée; Docker/Conda; \texttt{calibration.json}.
\item \textbf{Audit} : \texttt{report\_validation.json} contenant données agrégées, critères, verdicts; PDF auto-généré.
\end{itemize}

\subsection{Critères d'acceptation (exemple)}
\begin{itemize}[nosep]
\item Identification: macro-F1 $\ge 0.92$; ECE $\le 5\%$; TPR@FPR=5\% $\ge 90\%$ pour OOD.
\item Quantification: $R^2 \ge 0.95$ sur val. indépendante; RMSE $\le$ cible métier; biais $\approx 0$; LOD/LOQ respectées.
\item Métrologie: PV (Gage R\&R) $\le 10\%$; stabilité (EWMA) dans les bornes.
\end{itemize}

\subsection{Sorties \& formats}
\begin{itemize}[nosep]
\item \textbf{report\_validation.json} : DOE, métriques (avec IC), verdict pass/fail par critère.
\item \textbf{tables/} : tableaux normalisés (Gage R\&R, stabilité, robustesse).
\item \textbf{figures/} : cartes de stabilité, courbes Bland--Altman, distributions de métriques.
\end{itemize}

\subsection{Pseudo-code (génération rapport)}
\begin{lstlisting}[basicstyle=\ttfamily\small]
def build_validation_report(datasets, cfg):
    rr = compute_repeatability(datasets.rep)   # RSD, ANOVA components
    ri = compute_reproducibility(datasets.inter)
    stab = control_charts(datasets.stability)  # Shewhart/EWMA
    idm = id_metrics(datasets.id)              # F1/ECE/OOD
    qtm = quant_metrics(datasets.quant)        # RMSE/R2/bias/LOD/LOQ
    robust = hp_sensitivity_scan(datasets.hp_grid)
    verdict = apply_acceptance(rr, ri, stab, idm, qtm, robust, cfg.thresholds)
    write_json("artifacts/validation/report_validation.json",
               {"rr": rr, "ri": ri, "stability": stab,
                "id": idm, "quant": qtm, "robust": robust,
                "verdict": verdict, "criteria": cfg.thresholds})
    return verdict
\end{lstlisting}

\subsection{Tests \& DoD}
\begin{itemize}[nosep]
\item \textbf{Reproductibilité} : mêmes entrées $\Rightarrow$ mêmes rapports (hashs), $\pm$ float.
\item \textbf{Couverture} : IC cohérentes par bootstrap/ANOVA.
\item \textbf{Signalements} : alarmes cohérentes si dérive simulée; verdicts pass/fail corrects lorsque seuils franchis.
\item \textbf{DoD} : \texttt{report\_validation.json} \& figures/export générés; annexes tables présentes.
\end{itemize}
